% Einheit 3 – Nr. 40–49
\setcounter{aufgabe}{39}
\einheitenkopf[e3][3 Exponentialfunktion]{Einheit 3 von 5 · Exponentialfunktion aufstellen und auswerten}
\zweigzeile{Hier lernst du, eine Wachstumsgleichung $N(t) = N_0 \cdot q^t$ aufzustellen, zu deuten und damit Werte und Schwellen zu berechnen · neu in diesem Jahr · P10 · baut auf: Wachstumsfaktor und Tabelle (Einheit 2), Potenz mit dem Taschenrechner}

\begin{aufgabe}{Ich erkenne, welches Jahr der Schritt null ist.}
Unterstreiche das Startjahr und kreise das gesuchte Jahr ein. Rechne nichts.
\begin{teile}
\teil 2020 hatte ein Verein 150 Mitglieder. Die Zahl wächst jährlich um 5\,\%. Wie viele Mitglieder hat er 2025?
\teil Im Jahr 2018 lebten 900 Störche in einer Region, ihre Zahl nimmt jährlich um 2\,\% ab. Wie viele sind es im Jahr 2030?
\teil Wie groß ist im Jahr 2031 ein Guthaben, das 2024 bei 1500\,€ lag und seitdem jährlich mit 3\,\% verzinst wird?
\teil Ein Auto kostete 2022 neu 30\,000\,€. Welchen Wert hat es 2027, wenn es jährlich 15\,\% seines Werts verliert?
\end{teile}
\end{aufgabe}

\verfahren{Wachstumsgleichung aufstellen und deuten}

\begin{aufgabe}{Ich kann eine Wachstumsgleichung aufstellen.}
Gib den Anfangswert und den Wachstumsfaktor an und stelle die Gleichung auf.
\begin{teile}
\teil In einer Probe sind 100 Bakterien, ihre Zahl verdoppelt sich jede Stunde. \\ \feld{N_0} \quad \feld{q} \quad $N(t) = $ \leerfeld
\teil Auf einem Teich schwimmen 50 Wasserlinsen, ihre Zahl verdreifacht sich jede Woche. \\ \feld{N_0} \quad \feld{q} \quad $N(t) = $ \leerfeld
\teil Ein Guthaben von 400\,€ wächst jährlich um 10\,\%. \\ \feld{N_0} \quad \feld{q} \quad $N(t) = $ \leerfeld
\teil Eine Stadt mit 2500 Einwohnern wächst jährlich um 2\,\%. \\ \feld{N_0} \quad \feld{q} \quad $N(t) = $ \leerfeld
\teil Im Blut sind 90\,mg eines Wirkstoffs, stündlich werden 20\,\% abgebaut. \\ \feld{N_0} \quad \feld{q} \quad $N(t) = $ \leerfeld
\teil Ein Gletscher bedeckt 12\,km², die Fläche nimmt jährlich um 1,5\,\% ab. \\ \feld{N_0} \quad \feld{q} \quad $N(t) = $ \leerfeld
\end{teile}
\end{aufgabe}

\begin{aufgabe}{Ich kann die passende Gleichung unter mehreren auswählen.}
Kreuze die passende Gleichung an. Dabei ist $x$ die Zeit in Jahren und $y$ der Bestand.
\begin{teile}
\teil Ein Bestand von 50 Tieren verdoppelt sich jedes Jahr. \\ \kreuz{$y = 50 \cdot 2^x$}\quad\kreuz{$y = 2 \cdot 50^x$}\quad\kreuz{$y = 50 \cdot x^2$}\quad\kreuz{$y = 50 + 2x$}
\teil Ein Bestand von 600 Tieren nimmt jährlich um 5\,\% ab. \\ \kreuz{$y = 600 \cdot 1{,}05^x$}\quad\kreuz{$y = 600 \cdot 0{,}95^x$}\quad\kreuz{$y = 600 \cdot 0{,}5^x$}\quad\kreuz{$y = 0{,}95 \cdot 600^x$}
\teil Ein Kapital von 3000\,€ wächst jährlich um 1,2\,\%. (P10 2017 OS) \\ \kreuz{$y = 3000 \cdot 1{,}2^x$}\quad\kreuz{$y = 3000 \cdot 1{,}012^x$}\quad\kreuz{$y = 3000 \cdot 0{,}012^x$}\quad\kreuz{$y = 3000 + 1{,}012x$}
\end{teile}
\end{aufgabe}

\begin{aufgabe}{Ich kann die Teile einer Wachstumsgleichung im Sachzusammenhang deuten.}
Gib an, was die Zahlen und die Variable bedeuten.
\begin{teile}
\teil $N(t) = 300 \cdot 2^t$ beschreibt die Zahl der Hefezellen in einer Probe, $t$ in Stunden.
\teil $W(t) = 18\,000 \cdot 0{,}85^t$ beschreibt den Wert eines Autos in Euro, $t$ in Jahren.
\teil Erkläre, warum der Faktor 1,04 für eine Zunahme um 4\,\% steht.
\teil $f(x) = 50 \cdot 1{,}03^x$ gibt die Masse eines Fohlens in kg an, $x$ in Wochen nach der Geburt. Gib die Bedeutung von 50, von 1,03 und von $x$ an. (P10 2019 OS)
\end{teile}
\end{aufgabe}

\verfahren{Werte und Schwellenwerte berechnen}

\begin{aufgabe}{Ich kann mit einer Wachstumsgleichung den Wert nach einigen Schritten berechnen.}
Berechne den Wert.
\begin{teilezwei}
\tz $N(t) = 5 \cdot 2^t$; \ $N(3)$ \leerfeld & \tz $N(t) = 10 \cdot 3^t$; \ $N(2)$ \leerfeld \\
\tz $N(t) = 100 \cdot 1{,}1^t$; \ $N(2)$ \leerfeld & \tz $N(t) = 1000 \cdot 0{,}5^t$; \ $N(3)$ \leerfeld \\
\end{teilezwei}
\anweisung{Rechne mit dem Taschenrechner und runde erst am Ende sinnvoll.}
\begin{teile}
\teil $N(t) = 250 \cdot 1{,}08^t$; \ $N(10)$ \leerfeld
\teil Eine Stadt hatte im Jahr 2021 genau 4000 Einwohner, die Zahl wächst jährlich um 1,5\,\%. Wie viele Einwohner hat sie im Jahr 2029? \leerfeld
\teil Ein Auto kostet neu 24\,000\,€ und verliert jährlich 12\,\% seines Werts. Wie viel ist es nach fünf Jahren noch wert? \leerfeld[€]
\end{teile}
\end{aufgabe}

\begin{aufgabe}{Ich kann eine Behauptung über ein Wachstum mit einer Rechnung prüfen.}
Prüfe mit einer Rechnung und entscheide.
\begin{teile}
\teil $N(t) = 300 \cdot 1{,}5^t$. Behauptung: Nach vier Schritten ist der Bestand größer als 1500.
\teil Eine Gemeinde hat 5000 Einwohner und wächst jährlich um 2\,\%. Studie A sagt: „In zehn Jahren sind es 20\,\% mehr, also 6000.“ Studie B rechnet mit dem Wachstumsfaktor. Welche Studie sagt mehr voraus, und um wie viele Einwohner?
\teil Algen bedecken 30\,m² eines Sees, die Fläche wächst täglich um 25\,\%. Ein Zeitungsartikel behauptet: „Nach drei Wochen bedecken die Algen mehr als 3000\,m².“ Stelle eine Gleichung auf und prüfe die Behauptung. (P10 2025 OS)
\end{teile}
\end{aufgabe}

\begin{aufgabe}{Ich kann berechnen, wann ein Wert eine Schwelle überschreitet.}
Bestimme, nach wie vielen Schritten die Schwelle zum ersten Mal überschritten ist, und gib den Wert dann an.
\begin{teile}
\teil $N(t) = 10 \cdot 2^t$, Schwelle 100 \quad \feld{t} \feld{N}
\teil $N(t) = 3 \cdot 3^t$, Schwelle 200 \quad \feld{t} \feld{N}
\teil $N(t) = 500 \cdot 1{,}25^t$, Schwelle 1000 \quad \feld{t} \feld{N}
\end{teile}
\anweisung{Bestimme, nach wie vielen Schritten die Schwelle zum ersten Mal unterschritten ist.}
\begin{teile}
\teil $N(t) = 800 \cdot 0{,}7^t$, Schwelle 100 \quad \feld{t} \feld{N}
\end{teile}
\anweisung{Beantworte die Frage.}
\begin{teile}
\teil Ein Start-up hatte im Jahr 2022 einen Umsatz von 420\,000\,€ und steigert ihn jährlich um 15\,\%. In welchem Jahr übersteigt der Umsatz zum ersten Mal eine Million Euro? Gib den Umsatz in diesem Jahr an. (P10 2018 OS)
\end{teile}
\end{aufgabe}

\clearpage
\begin{aufgabe}{Ich finde den Fehler in einer Rechnung mit der Wachstumsgleichung.}
Finde den Fehler, benenne ihn und korrigiere die Rechnung.
\begin{teile}
\teil Ein Bestand von 600 Tieren wächst seit 2020 jährlich um 3\,\%. Jonas berechnet den Bestand im Jahr 2026:
\rechnung{N(t) &= 600 \cdot 1{,}03^t \\ N(2026) &= 600 \cdot 1{,}03^{2026}}
\teil Ein Guthaben von 1200\,€ wächst jährlich um 4\,\%. Sara berechnet das Guthaben nach fünf Jahren:
\rechnung{K &= 1200\,€ \cdot 1{,}04 \cdot 5 \\ K &= 6240\,€}
\teil Ein Bestand von 90 Tieren wächst jährlich um 4\,\%. Tim stellt die Gleichung auf:
\rechnung{N(t) &= 90 \cdot 4^t}
\end{teile}
\end{aufgabe}

\begin{aufgabe}{Ich kann begründen, warum in der Wachstumsgleichung eine Potenz steht.}
Begründe.
\begin{teile}
\teil Warum rechnet man bei einem Wachstum um 10\,\% in drei Jahren $500 \cdot 1{,}1^3$ und nicht $500 \cdot 1{,}1 \cdot 3$?
\teil Warum ist „dreimal nacheinander um 10\,\% erhöht“ mehr als „einmal um 30\,\% erhöht“?
\end{teile}
\end{aufgabe}

\begin{aufgabe}{Ich kann mit Wachstumsgleichungen zwei Entwicklungen vergleichen.}
Stadt A hat 40\,000 Einwohner und wächst jährlich um 1,5\,\%. Stadt B hat 46\,000 Einwohner und schrumpft jährlich um 1\,\%.
\begin{teile}
\teil Stelle für beide Städte eine Gleichung auf.
\teil Berechne beide Einwohnerzahlen nach fünf Jahren.
\teil Nach wie vielen Jahren hat Stadt A zum ersten Mal mehr Einwohner als Stadt B?
\end{teile}
\end{aufgabe}
